\documentclass[noauthor]{ximera}

\title{Week 3}
\author{Math 1193 Design Team}

\begin{document}
\begin{abstract}
    Solutions for Week 3.
\end{abstract}
\maketitle

\begin{problem}
Solve \(4x - 5y = 10\) for \(x\). Now solve it for \(y\).

\begin{explanation}
When we say solve an equation for \(x\), we mean isolate \(x\) on one side.

Here's some work:
\begin{align*}
    4x - 5y & = 10 \\
    4x & = 10 - 5y \\
    x & = \frac{10 - 5y}{4} \\
    x & = -\frac{5}{4}y + \frac{5}{2}
\end{align*}

The last line is not strictly necessary, but may be a nicer way to write it.

Solving for \(y\):
\begin{align*}
    4x - 5y & = 10 \\
    -5y & = 10 - 4x \\
    y & = \frac{10 -4x}{-5} \\
    y & = 2 + \frac{4}{5}x
\end{align*}
\end{explanation}
\end{problem}

\begin{problem}
Solve \(0.01x + 0.5y = 1\) for \(x\). Now solve it for \(y\).
Share with a group how you solved the equation. See if you can come up with more ways. Can you use fractions? Whole numbers?

\begin{explanation}
When we say solve an equation for \(x\), we mean isolate \(x\) on one side.

Here's some work:
\begin{align*}
    0.01x + 0.5y & = 1 \\
    0.01x & = 1 - 0.5y \\
    x & = \frac{1 - 0.5y}{0.01}
\end{align*}

We could also convert everything to fractions:
\begin{align*}
    \frac{x}{100} + \frac{y}{2} & = 1 \\
    100\left(\frac{x}{100} + \frac{y}{2}\right) & = 100 \cdot 1 \\
    x + 50y & = 100 \\
    x = 100 - 50y
\end{align*}

Decimals are similar.

Solving for \(y\):
\begin{align*}
    0.01x + 0.5y & = 1 \\
    0.5y & = 1 - 0.01x \\
    y & = \frac{1 -0.01x}{0.5} \\
    y & = \frac{1}{0.5} - \frac{0.01x}{0.5} \\
    y & = \frac{1}{1/2} - \frac{x/100}{1/2} \\
    y & = 1 \cdot \frac{2}{1} - \frac{x}{100}\cdot \frac{2}{1} \\
    y & = 2 - \frac{x}{50}
\end{align*}

You should be comfortable working with decimals, fractions, and converting between them.
\end{explanation}
\end{problem}

\begin{problem}
Let \(2a + b = 23\) and \(a = -b + 1\). 
\begin{enumerate}
    \item Rewrite the first equation in terms of \(b\). That is, rewrite it so no terms contain the variable \(a\).
    \item Rewrite the second equation in terms of \(a\). That is, rewrite it so no terms contain the variable \(b\).
\end{enumerate}

\begin{explanation}
    \begin{enumerate}
        \item Since \(a = -b + 1\), we can substitute into the first equation:
        \begin{align*}
            2(-b + 1) + b & = 23 \\
            -2b + 2 + b & = 23 \\
            -b + 2 & = 23 \\
            -b & = 21 \\
            b & = -21
        \end{align*}
        \item Now, from the first equation, \(b = 23 - 2a\). Substituting, we find that
        \begin{align*}
            a & = -(23 - 2a) + 1 \\
            a & = -23 + 2a + 1 \\
            a & = -22 + 2a \\
            a - 2a & = -22 \\
            -a & = -22 \\
            a & = 22
        \end{align*}
    \end{enumerate}
\end{explanation}
\end{problem}

\begin{problem}
\begin{enumerate}
    \item Graph the equations \(x + y = 2\) and \(x + y = 3\)
    without a calculator. Do these lines ever touch? Why or why not?
    Can you find an \(x\) and a \(y\) that satisfy both equations at 
    the same time? Why or why not? 
    \item Graph the equations \(x + y = 2\) and \(x - y = 3\)
    without a calculator. Do these lines ever touch? Why or why not?
    Can you find an \(x\) and a \(y\) that satisfy both equations at 
    the same time? Why or why not? 
    \item Is it possible for two lines in the 
    plane to touch more than once? Debate with a friend.
\end{enumerate}

\begin{explanation}
\begin{enumerate}
    \item Notice that the two equations are not in slope-intercept (\(y = mx + b\)) form, so one approach is to rearrange them into 
    that form: \(y = - x + 2\) and \(y = - x + 3\). Now you can see that these are lines with slope \(-1\) and \(y\)-intercepts 2 and 3, respectively.

    \begin{center}
    \includegraphics[width=10cm]{parallelLines.png}
    \end{center}

    From their graphs, these lines are parallel, so they do not ever touch.

    There are a few ways to think about the problem of finding an \(x\) and \(y\) that satisfy both equations at the same time. 
    
    Graphically, since the two lines don't intersect, there isn't an \((x, y)\) pair that lies on both of the graphs. But \((x, y)\) pairs
    are only on the graphs if they satisfy the equations and vice versa. Therefore, no \(x\) and \(y\) can satisfy both equations at the same time.

    You can also think about this by substituting. Since any \(y\) that satisfies both equations has to equal both \(-x + 2\) and \(-x + 3\) at the same time,
    by the transitive property of equality, \(-x + 2 = -x + 3\). Canceling the \(-x\), we arrive at the conclusion that \(2 = 3\), which
    can't be true. Therefore, no such \(y\) exists.

    \item Notice that the two equations are not in slope-intercept (\(y = mx + b\)) form, so one approach is to rearrange them into 
    that form: \(y = - x + 2\) and \(y = x - 3\). Now you can see that these are lines with slope \(-1\) and \(1\) and \(y\)-intercepts 2 and 3, respectively.

    \begin{center}
    \includegraphics[width=10cm]{perpLines.png}
    \end{center}

    From their graphs, these lines are not paralle, so they will touch. In fact, they're perpendicular, so they cross at a right angle.

    There are a few ways to think about the problem of finding an \(x\) and \(y\) that satisfy both equations at the same time. 
    
    Graphically, since the two lines intersect, there's an \((x, y)\) pair that lies on both of the graphs. But \((x, y)\) pairs
    are only on the graphs if they satisfy the equations and vice versa. Therefore, some \(x\) and \(y\) can satisfy both equations at the same time.

    You can also think about this by substituting. Since any \(y\) that satisfies both equations has to equal both \(-x + 2\) and \(x - 3\) at the same time,
    by the transitive property of equality, \(-x + 2 = x - 3\). Solving this equation gives us \(x = \frac{5}{2}\). Plugging into one of the original
    equations, we see that \(y = \frac{5}{2} - 3\), or \(y = -\frac{1}{2}\). This means that \(\left(\frac{5}{2}, -\frac{1}{2}\right)\) is
    an \((x, y)\) pair that satisfies both equations at the same time.

    \item 
\end{enumerate}
\end{explanation}
\end{problem}

\begin{problem}
Your second-favorite music artist is playing at a concert venue. There are two types of tickets. VIP seating costs \$100, and general admission costs \$40. 
Ticket sales totaled \$53,000. Your friend who works at the venue tells you they sold 50 VIP seats.

How many general admission tickets were sold?
\end{problem}


\begin{problem}
Leo has a solution that is 2\% bleach by volume (that is, for every 100mL, 2mL is bleach and 98mL is water). He also has one that is 5\% bleach by volume. 

How much of each solution (in mL) should Leo use if he wants to make 500mL of a 4\% bleach solution?

\begin{enumerate}
    \item Define variables that help you solve this 
    problem. Be specific about what quantities the variables
    represent. Compare with a friend and discuss why you made 
    the choices you made.
    \item Write equations that relate the variables you 
    have defined.
\end{enumerate}

%begin{problem}
%A theater charges \$50 per ticket for seats in Section A, 
%\$30 per ticket for seats in Section B, and \$20 per ticket 
%for seats in Section C. For one play, 4000 tickets were sold 
%for a total of \$120,000 in revenue. 

%\begin{enumerate}
%    \item If 1000 more tickets in Section B were sold 
%    than the other two tickets combined, how many tickets 
%    in each section were sold?
%    \item Define variables that help you solve this 
%    problem. Be specific about what quantities the variables
%    represent. Compare with a friend and discuss why you made 
%    the choices you made.
%    \item Write equations that relate the variables you 
%    have defined.
%\end{enumerate}

\begin{explanation}
Closely mirrors 1148 exam-type questions with a bit more scaffolding.
\end{explanation}
\end{problem}

\end{document}
